\section{Square norm families, exact parity obstructions, and a dynamic cover tower}
\label{qg-section}

This section gives a positive infinite family whose second norm is an actual
square, and a precise obstruction to making both norms squares up to the
ramified factor three. The proofs use explicit elliptic curves and descent.
A separate multiquadratic tower gives fixed-square-class finiteness for
successive mixed transforms. None of these conclusions supplies the unbounded
family required by the two-step radical criterion. The geometric arguments
are ordinary mathematics; any companion Lean arithmetic identities have
their narrower scope stated separately.

Put
\[
 N(x)=x^2+x+1,\qquad F(x)=x^4+3x^3+5x^2+3x+1.
\]
For coprime positive integers $a,b$, write $c=a+b$, $U=ab$ and
\[
 M_0=a^2+ab+b^2,\qquad
 M_1=U^2+UM_0+M_0^2
     =a^4+3a^3b+5a^2b^2+3ab^3+b^4.
\]

\subsection{An actual infinite family with square second norm}
\begin{lemma}[Explicit elliptic coordinates]\label{qg-birational}
The smooth projective curve $y^2=F(x)$ is birational over $\mathbb Q$ to
\begin{equation}\label{qg-elliptic-model}
 E:\quad Y^2=X^3-X^2-3X.
\end{equation}
On the open sets where the displayed denominators are nonzero, the maps are
\begin{align}
 u&=\frac{y-1-3x/2}{x^2},\quad X=2-2u,
 &Y&=\frac X2\bigl((X-4)x-3\bigr),\label{qg-forward}\\
 x&=\frac{2Y+3X}{X(X-4)},
 &y&=1+\frac32x+\left(1-\frac X2\right)x^2.
 \label{qg-inverse}
\end{align}
\end{lemma}
\begin{proof}
Substitute $y=1+3x/2+ux^2$ and divide the quartic equation by $x^2$:
\[
 (u^2-1)x^2+3(u-1)x+2u-\frac{11}4=0.
\]
The discriminant of this quadratic in $x$ is
$-8u^3+20u^2-10u-2$. Substitution $X=2-2u$ changes this to
$X^3-X^2-3X$, and completing the square gives \eqref{qg-forward}.
Equivalently, after $u=1-X/2$ the quartic difference is
\[
 \frac{x^2}4\bigl(X(X-4)x^2-6Xx-4X-3\bigr).
\]
Solving for $x$ gives \eqref{qg-inverse}; the two maps are inverse on their
open domains. The discriminant of $F$ is $117$, and the cubic in
\eqref{qg-elliptic-model} also has distinct roots, so their projective
normalizations are smooth genus-one curves.
\end{proof}

\begin{theorem}[Infinitely many positive square second norms]
\label{qg-square-family}
There are infinitely many coprime positive pairs $(a,b)$, with unbounded
$a+b$, for which $M_1$ is an integer square. Each such pair gives an actual
second Eisenstein extraction with exponent two and unit residual, hence
$\lambda_1=1/2$ for that extraction.
\end{theorem}
\begin{proof}
On $E$ take $P=(3,3)$. Exact group operations give
\[
 2P=(4,-6),\qquad 3P=(75,645),\qquad
 4P=\left(\frac{361}{144},\frac{2413}{1728}\right).
\]
The doubling formula is
\begin{equation}\label{qg-doubling}
 X(2Q)=\frac{(X(Q)^2+3)^2}{4X(Q)(X(Q)^2-X(Q)-3)}.
\end{equation}
If $s=v_2(X(Q))<0$, the numerator and denominator have valuations $4s$
and $2+3s$. Thus $v_2(X(2Q))=s-2$. Starting at $4P$, whose $X$ coordinate
has valuation $-4$, gives infinitely many distinct doubles. Therefore $P$
has infinite order without a database rank computation.

The real identity component contains infinity and the branch
$X\ge(1+\sqrt{13})/2$, hence contains $P$. It is a compact connected
one-dimensional Lie group, isomorphic to a circle by integration of an
invariant differential. The infinite-order point $P$ corresponds to an
irrational angle, so its integer multiples are dense. At $3P$, the inverse
map gives
\[
 x=\frac{101}{355}>0,\qquad y=-\frac{192529}{126025}.
\]
A neighborhood avoids the exceptional denominators and maps continuously
to $x>0$. Infinitely many distinct multiples lie there. Birationality and
the fact that each $x$ supports at most two $y$ give infinitely many
distinct positive rational ratios $x=a/b$ in lowest terms.
Now $(yb^2)^2=F(a,b)\in\mathbb Z$, and a rational number whose square is
integral is integral. Thus $M_1$ is an integer square. Infinitely many
reduced positive ratios have unbounded numerator plus denominator.

The first mixed output $(ab,M_0,c^2)$ is primitive and its norm $M_1$ has
no factor three. In its oriented Eisenstein factorization all split
exponents are therefore even. Unique factorization gives
$z_1=\text{unit}\cdot W^2$, so the specified unit-residual extraction has
$\lambda_1=1/2$.
\end{proof}
For example, $F(101,355)=192529^2$. Exponent two alone is much weaker than
the NT4 sufficient threshold; no stronger extraction or high ABC quality
is inferred. Exact searches of the existing repository found no prior use
of this pair or elliptic model in this connection. The elliptic curve itself
is familiar, and no assertion of literature novelty is made.

\subsection{Two explicit descent calculations}
We use Mordell--Weil finite generation and prime-to-$p$ torsion injectivity
at good reduction, as stated in the author-hosted J.~S.~Milne,
\emph{Elliptic Curves}, second edition, Chapter~IV, finite-basis theorem,
and Chapter~II, Corollary~4.2:
\url{https://www.jmilne.org/math/Books/EC2.pdf}.
The following descent maps and local exclusions are explicit, so no rank
or torsion assertion is taken from a database.

For integers $a,b$, let $E(a,b):y^2=x^3+ax^2+bx$, where
$b(a^2-4b)\ne0$, and let
$E'=E(-2a,a^2-4b)$. The two-isogeny is
\begin{equation}\label{qg-isogeny}
 \phi(x,y)=\left(x+a+\frac bx,\ y\left(1-\frac b{x^2}\right)\right).
\end{equation}
Applying the same formula to $E'$ and scaling by $1/4,1/8$ gives the dual
$\psi$, with $\psi\phi=[2]$. The square-class map
$\alpha:E(\mathbb Q)\to\mathbb Q^\times/\mathbb Q^{\times2}$ takes
the class of $x$ away from $O,(0,0)$ and satisfies
$\alpha(O)=1$, $\alpha((0,0))=[b]$. Its kernel is $\psi(E'(\mathbb Q))$.

Here the image criterion can be checked directly. The first coordinate in
\eqref{qg-isogeny} is $(y/x)^2$. Conversely, for a target $X=t^2\ne0$, the
equation $x^2+(a-X)x+b=0$ has discriminant
$(X-a)^2-4b=(Y/t)^2$. Choosing its root and the sign in $y=tx$ gives a
rational preimage. A target $(0,0)$ has a preimage exactly when the target's
$b$ coefficient is square, in agreement with its square-class value.
The homomorphism property follows from the product of the three $x$
coordinates cut out by a line $y=mx+n$ being $n^2$. For $n=0$, the product
of the other two roots is $b$, which is consistent with the assigned
value at $(0,0)$; vertical lines and tangencies give the remaining cases.
More explicitly, $T=(0,0)$ on $E$ lies in $\psi(E')$ exactly when $b$ is
square, and $O$ always lies in the image. The formula
$x(P+T)=b/x(P)$ verifies the class law through $T$, while vertical lines
give $\alpha(P)\alpha(-P)=1$.

Every image class is a signed squarefree divisor $d$ of $b$. A prime
outside $b$ cannot have odd positive valuation in $x$, and a negative
valuation in $x$ is even by comparing $y^2$ with the monic cubic.
Writing $x=dU^2/V^2$ and clearing denominators gives, for any such
nontrivial class, primitive integers satisfying
\begin{equation}\label{qg-covering-quartic}
 W^2=dU^4+aU^2V^2+(b/d)V^4,\qquad \gcd(U,V)=1.
\end{equation}
The covering coordinate $W$ is initially rational with integral square,
so it is an integer.

\begin{lemma}[Complete points on two small twists]\label{qg-small-twists}
The rational points of $E_1:y^2=x(x-1)(x+3)$ are exactly
\[
 O,\ (0,0),\ (1,0),\ (-3,0),\ (-1,\pm2),\ (3,\pm6).
\]
The rational points of $E_3:y^2=x(x-3)(x+9)$ are exactly
$O,(0,0),(3,0),(-9,0)$.
\end{lemma}
\begin{proof}
For $E_1$ one has $a=2,b=-3$, and its dual has $a'=-4,b'=16$.
The image of $\alpha$ consists of the four classes $1,-1,3,-3$, all
realized by the listed points. On the dual, $x\ge0$ because
$x^2-4x+16=(x-2)^2+12>0$. Thus the possible image classes are $1,2$.
Class two would force
\[
 W^2=2U^4-4U^2V^2+8V^4.
\]
In the primitive parity cases (odd,even), (odd,odd), (even,odd), the right
side is respectively $2,6,8$ modulo sixteen, all nonsquares. The dual
image is therefore trivial and $E_1'(\mathbb Q)=\phi(E_1(\mathbb Q))$.
Hence $\ker\alpha=\psi(E_1')=2E_1$ and
$|E_1(\mathbb Q)/2E_1(\mathbb Q)|=4$. There are four rational
two-torsion points; finite generation then gives rank zero.
Both good reductions at five and seven have eight points. The two
prime-to-$p$ torsion injections eliminate odd torsion and bound the
two-primary torsion by eight. All eight listed points occur, so the list
is complete.

For $E_3$, $a=6,b=-27$, and its dual has $a'=-12,b'=144$.
The four rational two-torsion points already realize the four classes
$1,-1,3,-3$ of $\alpha$. The dual has $x\ge0$, since
$x^2-12x+144=(x-6)^2+108>0$, and its possible classes are $1,2,3,6$.
For classes two and six the quartics in \eqref{qg-covering-quartic} are
\[
 2U^4-12U^2V^2+72V^4,\qquad
 6U^4-12U^2V^2+24V^4.
\]
Their residues for the same three parity cases are respectively
$(2,14,8)$ and $(6,2,8)$ modulo sixteen, excluding both.
For class three one would have
\[
 W^2=3U^4-12U^2V^2+48V^4.
\]
Writing $W=3W_0$ and dividing by three gives, modulo three,
$0=U^4-U^2V^2+V^4$. Its right side is one unless both $U,V$ are
divisible by three, contradicting primitivity. Thus the dual image is
again trivial, the quotient by doubling has order four, and the rank is
zero. The good reductions at five and seven have four and eight points,
so the torsion order divides four. Its four rational two-torsion points
exhaust the group.
\end{proof}

\begin{theorem}[The combined norm is neither a square nor three times one]
\label{qg-combined-obstruction}
For every positive rational $x$, neither $N(x)F(x)$ nor $N(x)F(x)/3$ is
a rational square. Consequently, for every primitive positive integer
seed $(a,b)$, neither $M_0M_1$ nor $M_0M_1/3$ is a rational square.
\end{theorem}
\begin{proof}
Suppose $N(x)F(x)=dh^2$ with $d\in\{1,3\}$. Set
\[
 u=x+x^{-1},\qquad v=2u+3,\qquad t=4h(x+1)/x^2.
\]
Since $N(x)/x=u+1$ and $F(x)/x^2=u^2+3u+3$,
\begin{equation}\label{qg-reciprocal-identity}
 dt^2=16(u+2)(u+1)(u^2+3u+3)
     =v^4+2v^2-3=(v^2-1)(v^2+3).
\end{equation}
Thus $(X,Y)=(v^2,vt)$ lies on $dY^2=X(X-1)(X+3)$.
For $d=1$ this is $E_1$; for $d=3$ the coordinates $(3X,9Y)$ lie on
$E_3$. Positivity gives $u\ge2,v\ge7$. The respective $x$ coordinates
are at least $49$ and $147$, contradicting Lemma~\ref{qg-small-twists}.
Finally $M_0M_1=b^6N(a/b)F(a/b)$, and $b^6$ is a square.
\end{proof}

\begin{corollary}[A precise even-exponent obstruction]
\label{qg-even-exponents}
In actual successive norm representations
\[
 M_0=3^eV_0Q_0^{g_0},\qquad M_1=V_1Q_1^{g_1},
 \qquad e\in\{0,1\},
\]
the exponents $g_0,g_1$ cannot both be even when both residual norms
$V_0,V_1$ are squares. In particular, two unit-residual even extractions
are impossible, even with arbitrarily moving roots.
\end{corollary}
\begin{proof}
Those assumptions would make $M_0M_1$ equal to $3^e$ times a square,
contrary to Theorem~\ref{qg-combined-obstruction}.
\end{proof}
This excludes a specific exact route, not arbitrary residuals, odd
extraction exponents, or the general finite lambda threshold.

\subsection{Fixed residual square classes and a genus-three curve}
\begin{theorem}[Fixed-class finiteness]\label{qg-fixed-classes}
For any fixed positive rational square classes $d_0,d_1$, only finitely
many coprime positive pairs $(a,b)$ satisfy
$M_0=d_0A^2$, $M_1=d_1B^2$ with rational $A,B$.
No effective height bound is asserted.
\end{theorem}
\begin{proof}
Take the smooth projective normalization with function field
\[
 \mathbb Q(x)\left(\sqrt{N(x)/d_0},\sqrt{F(x)/d_1}\right).
\]
The discriminants are $-3$ and $117$. The polynomials have no common root,
since $F\bmod N=x^2$ and $N(0)=1$; indeed their resultant is one.
Disjoint simple zeros prove that their geometric square classes are
independent. The cover is therefore geometrically connected of degree four.
There are six finite branch points. Each has two points of ramification
index two, contributing two to Riemann--Hurwitz. Infinity is unramified
because both degrees are even. Thus $2g-2=-8+12=4$, so $g=3$.
Faltings's theorem gives finitely many rational points on this fixed curve.
Every seed gives $x=a/b$ with square coordinates $A/b,B/b^2$, and a reduced
positive ratio determines its primitive pair uniquely.
\end{proof}
The external finiteness input is G.~Faltings,
\emph{Endlichkeitss\"atze f\"ur abelsche Variet\"aten \"uber Zahlk\"orpern},
Inventiones Mathematicae \textbf{73} (1983), 349--366, with the 1984 erratum;
the original article is indexed at \url{https://eudml.org/doc/143051}.

For two actual even extractions the square classes are
$([3^eV_0],[V_1])$. Along a family whose seed height tends to infinity,
this \emph{pair} must eventually leave each fixed finite product of
square-class sets. The quantifier is joint: it need not be one fixed
coordinate that escapes along the whole sequence. Moving square classes
and odd exponents remain possible.

\subsection{A dynamic tower with explicit genus}
Define
\begin{equation}\label{qg-polynomial-recursion}
 A_0=x,\ B_0=1,\quad
 A_{j+1}=A_jB_j,\quad B_{j+1}=A_j^2+A_jB_j+B_j^2,
 \quad f_j=B_{j+1}.
\end{equation}
Then $A_j/B_j$ is the $j$th iterate of $r(x)=x/(x^2+x+1)$, and
$f_0=N,f_1=F$.

\begin{theorem}[Dynamic multiquadratic covers]\label{qg-dynamic-tower}
Every $f_j$ is squarefree of degree $2^{j+1}$, and distinct $f_j$ have
disjoint root sets. For fixed nonzero rational $d_0,\ldots,d_{k-1}$, the
smooth projective cover with function field
\[
 \mathbb Q(x)\left(\sqrt{f_0/d_0},\ldots,
                         \sqrt{f_{k-1}/d_{k-1}}\right)
\]
is geometrically connected of degree $2^k$ and has genus
\begin{equation}\label{qg-tower-genus}
 G_k=1+2^{k-2}(2^{k+1}-6),\qquad k\ge1.
\end{equation}
In particular $G_1=0,G_2=3,G_3=21$. For $k\ge2$, fixing this vector
of square classes allows only finitely many primitive positive seeds for
the successive mixed transforms.
\end{theorem}
\begin{proof}
The degree-two map $r$ has exactly the critical points $1,-1$, since
$r'(x)=(1-x^2)/(x^2+x+1)^2$. The point $-1$ is fixed, and the forward
orbit of $1$ remains positive and finite. Neither orbit reaches infinity;
infinity maps to the fixed point zero. Thus infinity is not postcritical.
All points in $(r^{j+1})^{-1}(\infty)$ are unramified, giving exactly
$2^{j+1}$ distinct finite points.

The numerator and denominator in \eqref{qg-polynomial-recursion} are
coprime by induction: a common zero after one step would be a common zero
before it. For $j\ge1$ their degrees are $2^j-1,2^j$, with monic leading
coefficients. Hence these preimages are exactly the roots of $f_j$ and
$\deg f_j=2^{j+1}$. If two different levels shared a root, a positive
iterate of infinity would return to infinity, contrary to its orbit.

Any nonempty product of distinct $f_j$ has an odd valuation at a root of
one selected factor. Thus their geometric square classes are independent,
proving degree $2^k$ and geometric connectedness. There are
$2^{k+1}-2$ finite branch points, each with inertia group of order two and
total ramification contribution $2^{k-1}$. All degrees are even, so
infinity is unramified. Riemann--Hurwitz gives
\[
 2G_k-2=-2\cdot2^k+(2^{k+1}-2)2^{k-1},
\]
which is \eqref{qg-tower-genus}. Faltings gives finiteness for $k\ge2$.
Homogenizing $f_j(a/b)$ multiplies it by $b^{2^{j+1}}$, a square.
Consequently the classes are those of the actual successive norm integers.
\end{proof}
The tower gives explicit higher-genus objects and a fixed-class restriction.
It does not control the rational points uniformly over twists or settle the
moving-class two-step compression problem.
